\Article{28}
Since the formulas for $X$, $Y$, and $Z$ contain 71 terms in all, calculating directly from them is fairly laborious. Repeating the calculation for a large number of places would be all the more daunting, since without doing the same calculation twice one could hardly hope to be
\expandafter\def\csname HistoricalNoteText19\endcsname{Dr. \textsc{Goldschmidt} carried out the calculation for part of this auxiliary table.}
\SourceLeaf{157}safe from possible calculation errors. Little would be gained, either, by dropping all the terms whose coefficients are less than one unit, or even less than 10 units, since there would still be 65 terms left. But since the value of the work as a whole would remain uncertain if it were not tested against a considerable number of actual observations, I took the trouble to calculate an auxiliary table\HistoricalNote{19}. Using it greatly shortens and eases the work and thereby provides substantially better protection against calculation errors. Its arrangement is based on putting the values of the components in the following form
\[
\begin{aligned}
X={}&a^0+a'\cos(\lambda+A')+a''\cos(2\lambda+A'')\\
&+a'''\cos(3\lambda+A''')+a''''\cos(4\lambda+A'''')\\
Y={}&b'\cos(\lambda+B')+b''\cos(2\lambda+B'')\\
&+b'''\cos(3\lambda+B''')+b''''\cos(4\lambda+B'''')\\
Z={}&c^0+c'\cos(\lambda+C')+c''\cos(2\lambda+C'')\\
&+c'''\cos(3\lambda+C''')+c''''\cos(4\lambda+C'''')
\end{aligned}
\]
The first table contains the parts independent of $\lambda$ in $X$ and $Z$. The next four give the values of the auxiliary angles $A'$, $A''$ and so on, and the logarithms of $a'$, $a''$ and so on, all for each degree of latitude $\varphi=90^0-u$. [In the present printing, the table is appended to the end of this treatise in a slightly altered arrangement owing to the difference in format.]

The calculation for Göttingen may be given here as an example.

For the latitude $+51^0\,32'$, the tables give
\[
\begin{array}{r|r|r}
a^0=+500{,}8&&c^0=+1465{,}2\\
\log a'=2{,}28980&\log b'=2{,}18900&\log c'=2{,}20204\\
\log a''=1{,}79403&\log b''=2{,}03220&\log c''=2{,}12777\\
\log a'''=1{,}32522&\log b'''=1{,}46845&\log c'''=1{,}43199\\
\log a''''=0{,}59391&\log b''''=0{,}70016&\log c''''=0{,}59091\\
A'=249^0\,30'&B'=358^0\,24'&C'=105^0\,44'\\
A''=311\quad45&B''=64\quad50&C''=165\quad15\\
A'''=234\quad10&B'''=318\quad13&C'''=42\quad22\\
A''''=142\quad26&B''''=232\quad26&C''''=322\quad26
\end{array}
\]
\SourceLeaf{158}
and from these, with longitude $\lambda=9^0\,56\frac12'$, we get the contributions to
\[
\begin{array}{r|r|r}
X&Y&Z\\\hline
+500{,}8&&+1465{,}2\\
-35{,}71&+152{,}89&-68{,}99\\
+54{,}76&+9{,}92&-133{,}67\\
-2{,}21&+28{,}77&+8{,}27\\
-3{,}92&+0{,}19&+3{,}90\\\hline
X=+513{,}72&Y=+191{,}77&Z=+1274{,}71
\end{array}
\]
The rest of the calculation then gives
\[
\begin{aligned}
\delta&=+20^0\,28'\quad\log\omega=2{,}73907\\
i&=+66\quad43\\
\psi&=1387{,}6\quad\text{or in the usual unit}\\
\psi&=1{,}3876
\end{aligned}
\]

